\documentclass[11pt]{article}
\usepackage[margin=1in]{geometry}
\usepackage{amsmath,amssymb}
\usepackage{enumerate}
\usepackage{hyperref}
\usepackage{xcolor}
\usepackage{booktabs}
\usepackage{array}

%% ============================================================
%% Severity macros (replaces action macros for a code audit)
%% ============================================================
\newcommand{\critical}{\textcolor{red}{\textbf{[CRITICAL]}}}
\newcommand{\high}{\textcolor{orange}{\textbf{[HIGH]}}}
\newcommand{\medium}{\textcolor{blue}{\textbf{[MEDIUM]}}}
\newcommand{\low}{\textcolor{gray}{\textbf{[LOW]}}}
\newcommand{\confirmed}{\textcolor{teal}{\textbf{[CONFIRMED MATCH]}}}

\title{Math + Code Audit Report\\[6pt]
\large Gross \& Chivers (2025) ``NIMBYism and the Housing Supply''}
\author{}
\date{February 2026}

\begin{document}
\maketitle

\noindent This document reports the findings of a static code audit of
\texttt{Code/steadystate/} (MATLAB) and \texttt{Code/data/} (Stata)
against the paper's stated model, equations, and calibration table.
Seven MATLAB files and six Stata do-files were read in full.
No compile or run tests were executed.

\bigskip
\noindent\textcolor{orange}{\textbf{Scope note.}} \textit{%
\texttt{Code/Codes\_ABB/} (quantitative results, impulse responses,
calibration routines) was not covered in this pass and should be
audited separately. The \texttt{Code/steadystate/} folder is the
likely production codebase; see Outstanding Item~O1.}

\bigskip
\noindent\textbf{Summary of findings.}

\begin{center}
\begin{tabular}{lrrr}
\toprule
Severity & MATLAB & Stata & Total \\
\midrule
Critical & 4 & 2 & 6 \\
High     & 6 & 3 & 9 \\
Medium   & 3 & 3 & 6 \\
Low      & 2 & 2 & 4 \\
\midrule
Total    & 15 & 10 & 25 \\
\bottomrule
\end{tabular}
\end{center}

%% ============================================================
\section*{High-Level Math Audit (Paper First)}
%% ============================================================

\noindent This section checks internal model logic and equation coherence in the paper
before the code crosswalk. Source checked:
\texttt{Gross and Chivers (2025) NIMBYism and the Housing Supply.lyx}.

\begin{center}
\begin{tabular}{lll}
\toprule
ID & Severity & Description \\
\midrule
HLM1 & High & Median-voter regularity conditions not fully justified \\
HLM2 & Medium & Mixed marginal and finite-difference voting objects \\
HLM3 & Medium & Time/notation switching for prices and assets \\
HLM4 & Medium & Calibration prose mismatches listed parameter vector \\
HLM5 & Low & Math-adjacent typos in exposition \\
\bottomrule
\end{tabular}
\end{center}

\subsection*{HLM1: Median-voter regularity conditions}

\high{} The paper uses the equilibrium condition
\[
  \int g(\mathbf{s};p_t^h)\,\Theta(\mathbf{s};p^h)\,ds = 0.5,
\]
while also noting that household preferences over house prices are not monotonic for
some households. That can violate single-peakedness assumptions typically used to justify
a one-dimensional median-voter equilibrium.

\noindent\textbf{Inference.} This is not automatically a model error, but the paper should
state whether equilibrium is claimed as a numerical root condition only, or whether
existence/uniqueness follows from an explicit monotonicity assumption.

\subsection*{HLM2: Mixed marginal and finite-difference objects}

\medium{} The narrative refers to marginal effects (\(\partial V/\partial p\)), while
the formal voting rule uses:
\[
  V_j(\mathbf{s};p^h+\epsilon^p)-V_j(\mathbf{s};p^h)+\sigma \ge 0.
\]
These are equivalent only under a local-limit interpretation (\(\epsilon^p \to 0\)) or
additional smoothness assumptions.

\noindent\textbf{Requested clarification.} State whether \(\epsilon^p\) is an
infinitesimal perturbation used numerically, or a finite policy step.

\subsection*{HLM3: Time/notation consistency}

\medium{} The paper alternates between \(p_t\) and \(p_t^h\) for house price, and uses
both static and recursive constraint forms (\(a_t\) and \(a_{t+1}\) contexts). This likely
reflects exposition choices rather than a structural inconsistency, but it obscures timing.

\noindent\textbf{Requested clarification.} Add one compact timing-convention paragraph
(start-of-period states, controls, end-of-period transitions), and normalize price notation.

\subsection*{HLM4: Calibration prose vs parameter vector}

\medium{} The calibration sentence describes a narrower parameter class, but the listed
calibrated vector includes \(\beta, \sigma, \zeta, \bar{w}, \psi, \phi^{rent}\),
which spans preferences, voting, bequests, and rental-sector terms.

\noindent\textbf{Requested clarification.} Rewrite the sentence to classify all calibrated
parameters by block.

\subsection*{HLM5: Minor math-adjacent typos}

\low{} There are low-risk typographic issues in formal section headers and nearby prose
(e.g., ``Medan Voter Theorem''). These do not alter equations but reduce presentation quality.

\subsection*{High-Level Math Verdict}

No clear algebraic sign error was found in the core household Bellman setup, tenure-choice
decomposition, or rental zero-profit equation in this high-level pass. The main outstanding
math risk is the equilibrium-voting regularity argument (HLM1), followed by notation/timing
clarity (HLM2--HLM4).

%% ============================================================
\section*{Critical Issues (M1--M4; S1--S2)}
%% ============================================================

\subsection*{M1: Housing weight in utility function}

\critical{} The paper (Eq.~2) specifies a CES period utility function:
\[
  u(c_t, h_t) = \frac{\bigl(c_t^{1-\zeta}\,h_t^{\zeta}\bigr)^{1-\gamma}}{1-\gamma},
  \qquad \zeta = 0.75,
\]
so housing receives weight $\zeta = 0.75$ and consumption receives $1-\zeta = 0.25$.

The primary solver \texttt{SolveSS.m} (line~15) and \texttt{SolveSS\_iter.m} (line~14) implement:
\begin{verbatim}
omega = 0.5
value = (C^omega * (H + theta_r*R)^(1-omega))^(1-eta) / (1-eta)
\end{verbatim}
Here \texttt{omega} is the \emph{consumption} weight, so the housing weight is $1-\omega = 0.50$. The correct code translation of $\zeta = 0.75$ would be \texttt{omega = 0.25} (not 0.50). Both the parameter value and the role of the symbol are reversed relative to the paper.

\medskip
\noindent\textbf{Impact.} The consumption--housing substitution ratio is entirely determined by $\zeta$. A value of 0.50 rather than 0.75 substantially changes optimal housing demand, ownership rates, and the price--income ratio---all calibration targets.

\noindent\textbf{Decision needed.} Confirm the calibrated value of $\zeta$. If $\zeta=0.75$ is correct, set \texttt{omega = 1 - 0.75 = 0.25}.

\subsection*{M2: Discount factor $\beta$}

\critical{} The paper's calibration table states $\beta = 0.978$.

\begin{center}
\begin{tabular}{lll}
\toprule
File & Line & Value \\
\midrule
\texttt{SolveSS.m}          & 17 & \texttt{beta = 0.80} \\
\texttt{SolveSS\_function.m} & 16 & \texttt{beta = 0.98} \\
\texttt{SolveSS\_iter.m}    & 16 & \texttt{beta = 0.80} \\
\texttt{SolveSS\_Loop.m}    & 20 & \texttt{beta = 0.80} \\
\bottomrule
\end{tabular}
\end{center}

Only the optimisation wrapper \texttt{SolveSS\_function.m} is close to the calibrated value.
The remaining three files---including the main standalone solver---use $\beta = 0.80$,
which implies implausibly low patience and would dramatically compress life-cycle saving
and housing accumulation.

\noindent\textbf{Decision needed.} Identify the production file and update it to $\beta = 0.978$.

\subsection*{M3: Non-financial preference parameter $\sigma$ absent from code}

\critical{} The paper defines the individual voting indicator (Eq.~15) as:
\[
  \Theta(s;\,p^h) = \mathbf{1}\!\left\{V_j(s;\,p^h+\varepsilon) - V_j(s;\,p^h) + \sigma \geq 0\right\},
  \qquad \sigma = -2.4.
\]
The parameter $\sigma$ captures non-financial motives for voting (congestion, ideology)
and is calibrated to shift approximately 15\% of voters toward supply.

No variable named \texttt{sigma} or equivalent with value $-2.4$ appears in any
\texttt{Code/steadystate/} file. The voting block in \texttt{SolveSS.m}
(lines~393--474) computes:
\begin{verbatim}
vote = sign(valuefunction_dp - valuefunction) * population_density
\end{verbatim}
This captures only the financial component; the $\sigma$ shift is absent entirely.

\noindent\textbf{Impact.} Without $\sigma$, the voting equilibrium is determined
solely by financial incentives and will produce a different equilibrium price from the paper.

\noindent\textbf{Decision needed.} Confirm whether $\sigma$ was intentionally dropped
(e.g.\ absorbed into a price normalisation) or is missing. If missing, add as an additive
term in the voting threshold.

\subsection*{M4: Rental price hardcoded; zero-profit condition absent}

\critical{} The paper derives the rental price from a firm zero-profit condition (Eq.~12):
\[
  p_t^{\mathrm{rent}} = \varphi^{\mathrm{rent}} + r_t^a \cdot p_t^h - \Delta p_{t+1}^h,
\]
where $\varphi^{\mathrm{rent}} = 0.05$, $r_t^a = 0.02$, and $\Delta p_{t+1}^h$
is expected house price appreciation.

\texttt{SolveSS.m} (line~13) instead sets:
\begin{verbatim}
r_price = 0.09
\end{verbatim}
This is a fixed constant, not derived from the equilibrium condition. No file in
\texttt{Code/steadystate/} computes \texttt{r\_price} as a function of
\texttt{a\_price}, \texttt{rbPos}, or expected appreciation.

\noindent\textbf{Impact.} The rental price does not respond to changes in the house price
or interest rate, breaking the arbitrage between the rental and ownership markets.

\noindent\textbf{Decision needed.} Either implement Eq.~12, or document that 0.09 is
a pre-computed equilibrium value at the baseline calibration and explain why dynamic
updating is unnecessary.

\subsection*{S1: All Stata merges drop \texttt{\_merge} without validation}

\critical{} Every merge in the data pipeline immediately discards the match
indicator without tabulating match rates. Examples from \texttt{3 birthrates.do}:
\begin{verbatim}
merge m:1 statefips year using birthrates_updated   // line 25
drop _merge                                          // line 26

merge m:m metarea using weights, keepusing(sumcity) // line 41
drop _merge                                          // line 42
\end{verbatim}
The same pattern appears in files 5 and 6. There is no \texttt{tab \_merge},
\texttt{assert \_merge == 3}, or any check for unmatched observations.

\noindent\textbf{Impact.} Any silently dropped metropolitan areas or years will not
appear in the regression sample without re-running and inserting diagnostics.
This is a replication risk.

\noindent\textbf{Fix.} Add \texttt{tab \_merge} before every \texttt{drop \_merge}.
For key merges, add \texttt{assert \_merge == 3} or document the expected match rate.

\subsection*{S2: Impossible logical condition in \texttt{2 citystateweights.do}}

\critical{} The manual override block (around line~1234) contains at minimum one
condition that is logically impossible:
\begin{verbatim}
drop if metarea == 60 & bpl == 13 & bpl == 45
\end{verbatim}
A variable cannot simultaneously equal 13 and 45. This evaluates to false for all
observations and silently does nothing. The probable intent was:
\begin{verbatim}
drop if metarea == 60 & (bpl == 13 | bpl == 45)
\end{verbatim}

\noindent\textbf{Impact.} The intended observations are not dropped, potentially leaving
outlier or erroneous weight cells in the shift-share instrument. This affects weighted
birth rate construction and all downstream regressions.

\noindent\textbf{Fix.} Review the full 600-line block for similar errors and replace
\texttt{\&} with \texttt{|} wherever \texttt{bpl} conditions are clearly alternatives.

%% ============================================================
\section*{High Priority Issues (M5--M10; S3--S5)}
%% ============================================================

\subsection*{M5: Retirement age}

\high{} The paper states working life $J^{\mathrm{ret}} = 41$ periods (ages 25--65).
\texttt{SolveSS.m} (line~34) sets \texttt{agepension = 60}, implying retirement at
age~60 and only 35 working periods---five years shorter than the paper. This affects
the earnings profile, savings accumulation, and the age at which households transition
from renting to owning.

\subsection*{M6: Risk-free savings rate}

\high{} The calibration table states $r^a = 0.02$ (2\%). \texttt{SolveSS.m}
(line~10) sets \texttt{rbPos = 0.03} (3\%). This affects all saving decisions,
the rental price via Eq.~12, and the mortgage spread. With the paper's spread of
3~percentage points, a savings rate of 3\% implies a mortgage rate of 6\%, not 5\%.

\subsection*{M7: Bequest parameters}

\high{} The paper (Eq.~4 and calibration table) specifies:
\[
  v(w) = \psi\,\frac{(w + \bar{w})^{1-\gamma}}{1-\gamma}, \qquad \psi = 10,\quad \bar{w} = 150.
\]
\texttt{SolveSS.m} (lines~23--24) sets \texttt{bequestweight1 = 0.72} and
\texttt{bequestweight2 = 0.72}, and implements:
\begin{verbatim}
bequestweight1 * (1 + Bequest/bequestweight2)^(1-eta)
\end{verbatim}
Neither value is remotely close to $\psi = 10$ or $\bar{w} = 150$. The normalisation
$(1 + B/\bar{w})$ could be algebraically equivalent to $(w + \bar{w})^{1-\gamma}$
up to a scaling constant, but \texttt{bequestweight2} $= 0.72 \neq 150$.

\noindent\textbf{Decision needed.} Provide the algebraic equivalence mapping both
parameter sets, or correct the values.

\subsection*{M8: Transaction cost inconsistent across files}

\high{} The paper states $F^{\mathrm{sell}} = 0.07$ (7\%).

\begin{center}
\begin{tabular}{lll}
\toprule
File & Line & \texttt{ka} \\
\midrule
\texttt{SolveSS.m}          & 22 & 0.07 \checkmark \\
\texttt{SolveSS\_function.m} & 18 & 0.06 \\
\texttt{SolveSS\_Loop.m}    & 22 & 0.09 \\
\bottomrule
\end{tabular}
\end{center}

Two of three function files use a different value from the paper.
$F^{\mathrm{sell}}$ enters the budget constraint for adjusting owners (Eqs.~7 and~10)
and governs the margin between adjusting and non-adjusting tenure.

\subsection*{M9: Rental discount $\theta^{\mathrm{rent}}$ inconsistent across files}

\high{} The calibration table (following Mian et al.\ 2017) states $\theta^{\mathrm{rent}} = 0.90$.

\begin{center}
\begin{tabular}{lll}
\toprule
File & Line & \texttt{theta\_r} \\
\midrule
\texttt{SolveSS.m}          & 16 & 0.90 \checkmark \\
\texttt{SolveSS\_function.m} & 15 & 0.85 \\
\texttt{SolveSS\_Loop.m}    & 19 & 0.85 \\
\bottomrule
\end{tabular}
\end{center}

$\theta^{\mathrm{rent}}$ governs the utility discount for renting vs.\ owning and
directly shifts the rent--own margin.

\subsection*{M10: Utility form inconsistent across solver files}

\high{} The paper specifies CRRA utility with $\gamma = 2$ (Eq.~2). The files differ:

\begin{center}
\begin{tabular}{ll}
\toprule
File & Utility form \\
\midrule
\texttt{SolveSS.m}           & CRRA, $\eta = 2$ \\
\texttt{SolveSS\_iter.m}     & CRRA, $\eta = 2$ \\
\texttt{SolveSS\_function.m} & Log: $\log c + s_h \log(h + \theta_r R)$ \\
\texttt{SolveSS\_Loop.m}     & Log \\
\bottomrule
\end{tabular}
\end{center}

Log utility is the limiting case $\gamma \to 1$ and constitutes a different model.
If \texttt{SolveSS\_Loop.m} is the grid-search production solver, the equilibrium
prices it generates will not correspond to the CRRA model described in the paper.

\subsection*{S3: 600$+$ lines of manual \texttt{replace share = 0} with no documentation}

\high{} \texttt{2 citystateweights.do} lines~615--1234 contain over 600 statements of the form:
\begin{verbatim}
replace share = 0 if metarea == X & bpl == Y
\end{verbatim}
No comment explains the selection criterion. The logic is hand-coded, not
rule-based (e.g.\ \texttt{replace share = 0 if count < 5}).

\noindent\textbf{Impact.} The shift-share instrument is constructed from these weights.
Undocumented exclusions affect instrument validity and cannot be assessed by a reader
or replicator.

\noindent\textbf{Recommendation.} Replace the manual block with a programmatic rule
and document the threshold. If individual overrides are genuinely case-specific, add
a comment to each.

\subsection*{S4: Hardcoded machine-specific file paths}

\high{} Multiple files contain paths tied to specific machines:

\begin{center}
\begin{tabular}{lll}
\toprule
File & Line & Path \\
\midrule
\texttt{1 tenurechoices.do} & 1 & \texttt{/Users/igro0002/Dropbox/\ldots} \\
\texttt{4 importexcel .do} & 2, 13 & \texttt{/Users/igro0002/Dropbox/\ldots} \\
\texttt{merge data.do} & 6--7 & \texttt{C:\textbackslash{}Users\textbackslash{}Dave\_\textbackslash{}\ldots} \\
\bottomrule
\end{tabular}
\end{center}

The mixture of macOS and Windows paths indicates the pipeline was developed
across machines without standardisation. No collaborator or replicator can run
the code without modifying these lines.

\noindent\textbf{Fix.} Replace all hardcoded paths with a single global macro
at the top of each file, e.g.:
\begin{verbatim}
global datadir "C:\Users\Dave_\Dropbox\Zac and David\Data"
use "$datadir/usa_00009.dta"
\end{verbatim}

\subsection*{S5: IV regressions without instrument diagnostics}

\high{} \texttt{5 combinepermits.do} (lines~103--118) runs 50~IV regression
specifications in nested loops. No first-stage $F$-statistics are reported and
no Kleibergen--Paap or Cragg--Donald statistics are computed. The outlier
exclusion \texttt{abs(totalunits[h]growth) < 2} is not documented.

\noindent\textbf{Recommendation.} Add \texttt{estat firststage} to each IV loop
and export first-stage $F$-statistics alongside point estimates in
\texttt{Regressions.csv}.

%% ============================================================
\section*{Medium Priority Issues (M11--M13; S6--S8)}
%% ============================================================

\subsection*{M11: Housing grid inconsistent across files}

\medium{} The paper does not report the grid size used for published results.
The files vary substantially:

\begin{center}
\begin{tabular}{llll}
\toprule
File & $J$ & \texttt{housingmax} & Liquid grid $I$ \\
\midrule
\texttt{SolveSS.m}           & 40 & 14 & 50 \\
\texttt{SolveSS\_function.m} & 10 & 15 & 50 \\
\texttt{SolveSS\_iter.m}     & 20 & 25 & 50 \\
\texttt{SolveSS\_Loop.m}     & 20 & 30 & 60 \\
\bottomrule
\end{tabular}
\end{center}

A grid of $J = 10$ is very coarse and likely insufficiently accurate for value function
approximation. The varying \texttt{housingmax} changes which housing quantities are
accessible to households. Document which grid was used for the published results.

\subsection*{M12: \texttt{agemax = 90} in \texttt{SolveSS\_function.m}}

\medium{} The paper states agents live from age~25 to age~80 ($J = 56$ periods).
\texttt{SolveSS\_function.m} (line~26) sets \texttt{agemax = 90}, extending the
model life by 10~years. With \texttt{dage = 5} in that file (coarser age grid),
the terminal condition and bequest calculation apply at a different age than stated
in the paper.

\subsection*{M13: Notation reversal between paper and code}

\medium{} The paper uses $\zeta$ as the \emph{housing} weight ($\zeta = 0.75$,
so housing receives 75\%). The code uses \texttt{omega} as the
\emph{consumption} weight. The correct translation is \texttt{omega = 1 - zeta = 0.25}.
The code instead uses \texttt{omega = 0.50} (see~M1). The reversal in naming
conventions creates a risk of silent errors whenever the utility function is
re-parameterised.

\subsection*{S6: Lag construction duplicated between files~3 and~5}

\medium{} \texttt{3 birthrates.do} constructs lags \texttt{l1bw}--\texttt{l50bw}
(lines~55--119). \texttt{5 combinepermits.do} reconstructs \texttt{l1bw}--\texttt{l65bw}
(lines~13--77). If the panel changes between files due to the crosswalk merge,
lags may be computed over a different \texttt{tsset} structure. Any change to the
lag range in file~3 must be replicated manually in file~5.

\subsection*{S7: \texttt{statefips / 100} conversion}

\medium{} \texttt{3 birthrates.do} (line~21) converts IPUMS birth-state codes to
FIPS codes via \texttt{replace statefips = statefips / 100}. This arithmetic
conversion may fail for codes that are not round multiples of 100 (e.g.\
District of Columbia at IPUMS code~98 $\to$ 0.98, not a valid FIPS code).

\noindent\textbf{Recommendation.} Use the IPUMS-provided state crosswalk rather
than division. Verify all resulting values are valid FIPS codes before merging.

\subsection*{S8: Predicted variable created but unused}

\medium{} \texttt{5 combinepermits.do} (line~99) creates \texttt{own\_predicted}
after a fixed-effects regression but this variable is never used in any subsequent
command and is not saved to output. Remove or document its purpose.

%% ============================================================
\section*{Low Priority Issues (M14--M15; S9--S10)}
%% ============================================================

\subsection*{M14: \texttt{IncomeProcess.m} is an incomplete fragment}

\low{} \texttt{IncomeProcess.m} ends at line~104, mid-construction of lower-diagonal
arrays for an income transition matrix. There is no return statement, no assembled
sparse matrix, and no output. The file cannot be called as a standalone function.
It appears to be a copy-paste fragment that was never completed.

\subsection*{M15: Unused variable \texttt{sigma} in \texttt{SolveSS.m}}

\low{} \texttt{SolveSS.m} (line~18) declares \texttt{sigma = 1.5} but this variable
is never referenced again. The name coincides with the paper's non-financial preference
parameter $\sigma = -2.4$ (see~M3), but the value differs and it is unused. Either
remove or rename to avoid confusion.

\subsection*{S9: Commented-out Wharton merge in \texttt{merge data.do}}

\low{} Lines~28--32 of \texttt{merge data.do} contain a commented-out merge with a
Wharton land-use regulation dataset (WRLURI), listed in the paper as a potential
instrument for housing supply. No explanation is given for why it was removed.
Document the decision.

\subsection*{S10: Age bin boundary overlap}

\low{} \texttt{1 tenurechoices.do} creates \texttt{age20\_30} for ages 20--30
and \texttt{age30\_40} for ages 30--40. Age~30 falls in both bins. This may be
intentional but can cause double-counting in regression or summary tables.
Confirm and document.

%% ============================================================
\section*{Confirmed Matches}
%% ============================================================

The following items agree between the paper's calibration table and the primary
solver \texttt{SolveSS.m}:

\begin{itemize}
  \item \confirmed{} \textbf{LTV ratio $m$}: Paper $m = 0.90$ $\to$ \texttt{CC = 0.90} (line~26).
  \item \confirmed{} \textbf{Risk aversion $\gamma$}: Paper $\gamma = 2$ $\to$ \texttt{eta = 2}
    (\texttt{SolveSS.m}; \texttt{SolveSS\_iter.m}).
  \item \confirmed{} \textbf{Transaction cost (primary file)}: Paper $F^{\mathrm{sell}} = 0.07$
    $\to$ \texttt{ka = 0.07} (\texttt{SolveSS.m} line~22).
  \item \confirmed{} \textbf{Rental discount (primary file)}: Paper $\theta^{\mathrm{rent}} = 0.90$
    $\to$ \texttt{theta\_r = 0.90} (\texttt{SolveSS.m} line~16).
  \item \confirmed{} \textbf{Life span}: Paper $J = 56$ (ages 25--80) $\to$
    \texttt{agemin = 25}, \texttt{agemax = 80}, \texttt{dage = 1} (\texttt{SolveSS.m}).
  \item \confirmed{} \textbf{Borrowing constraint form}: Paper $a_t \geq -m \cdot h_t \cdot p_t^h$
    $\to$ \texttt{-b > a*P\_h*CC} (\texttt{SolveSS.m} lines~144--147).
  \item \confirmed{} \textbf{Backward induction structure}: Paper solution algorithm
    (backward from terminal age, forward distribution simulation) $\to$
    \texttt{SolveSS.m} backward loop from \texttt{agemax} to \texttt{agemin} followed
    by forward pass.
  \item \confirmed{} \textbf{Voting direction}: Paper votes to restrict supply when
    value rises with price $\to$ \texttt{sign(valuefunction\_dp - valuefunction)}
    (\texttt{SolveSS.m} lines~393--474). Note: financial component only; $\sigma$ term absent (see~M3).
\end{itemize}

%% ============================================================
\section*{Referee Recommendation (Code and Replication)}
%% ============================================================

\noindent\textbf{Recommendation:} \textbf{Major code revision} before treating the
current repository as a replication-ready implementation of the published model.

\medskip
\noindent This recommendation is based on observed implementation mismatches
(M1--M4), parameter inconsistencies across solver files (M5--M10), and missing
data-pipeline diagnostics (S1--S5). As an inference, if the published quantitative
results were generated from a different private branch, the immediate requirement is
to document the production-file mapping and provide a clean reproducible pipeline in
this repository.

\medskip
\noindent\textbf{Minimum acceptance checklist for replication release:}
\begin{itemize}
  \item Identify and document the exact production solver file(s) and calibration
    used for published tables/figures.
  \item Reconcile utility, discounting, voting, and rental-price equations with
    the paper, or explicitly document approved deviations.
  \item Add merge diagnostics and path standardisation to Stata scripts and re-run
    to regenerate key regression outputs.
  \item Publish a short \texttt{README\_code.md} mapping each paper equation block
    to code locations.
  \item Run and archive one end-to-end reproducibility test from raw inputs to
    final outputs.
\end{itemize}

%% ============================================================
\section*{Outstanding Items}
%% ============================================================

\begin{enumerate}
  \item \textbf{O1 (production file)}: \texttt{Code/steadystate/} is the
    likely production codebase, but \texttt{SolveSS.m} contains parameter values
    inconsistent with the calibration table (M1--M9). Verify whether
    \texttt{SolveSS.m} is the production file or a development variant;
    document in a \texttt{README\_code.md}. \textit{Owner: Authors.}

  \item \textbf{O2 (housing weight)}: Is $\zeta = 0.75$ (housing weight) or
    $\omega = 0.50$ (consumption weight) the calibrated value? Reconcile code
    against calibration table. \textit{Owner: Authors.}

  \item \textbf{O3 ($\sigma$ absent from code)}: Was $\sigma = -2.4$ intentionally
    dropped? If not, add as an additive term in the voting indicator before the
    median voter calculation. \textit{Owner: Authors.}

  \item \textbf{O4 (rental price)}: Is \texttt{r\_price = 0.09} a pre-computed
    equilibrium value or an error? If pre-computed, document the derivation.
    If an error, implement Eq.~12. \textit{Owner: Authors.}

  \item \textbf{O5 (retirement age)}: Retirement age 60 (code) or 65 (paper)?
    Update whichever is wrong. \textit{Owner: Authors.}

  \item \textbf{O6 (bequest parameters)}: Provide algebraic equivalence between
    $(\psi = 10,\, \bar{w} = 150)$ and $(\texttt{bequestweight1} = 0.72,\,
    \texttt{bequestweight2} = 0.72)$, or correct the code. \textit{Owner: Authors.}

  \item \textbf{O7 (\texttt{Codes\_ABB/} audit)}: The quantitative results,
    impulse responses, and calibration routines in \texttt{Code/Codes\_ABB/}
    were not covered in this pass. Schedule a separate audit before resubmission.
    \textit{Owner: Auditor.}

  \item \textbf{O8 (impossible bpl condition)}: Review and fix the manual override
    block in \texttt{2 citystateweights.do}; re-run the weights and all downstream
    files. \textit{Owner: Code author.}

  \item \textbf{O9 (merge diagnostics)}: Add \texttt{tab \_merge} before every
    \texttt{drop \_merge} across the Stata pipeline.
    \textit{Owner: Code author.}

  \item \textbf{O10 (share override documentation)}: Replace or document the
    600-line manual \texttt{replace share = 0} block with a programmatic rule
    or per-line comments. \textit{Owner: Code author.}
\end{enumerate}

\end{document}
